\documentclass[10pt, a4paper]{article}
% Page Layout
\usepackage[top=0.8cm, bottom=0.75cm, left=1cm, right=1cm]{geometry}
\usepackage{setspace} % Adjust line spacing with \onehalfspacing, \doublespacing, \setstretch{1.25}
% Language and Encoding
\usepackage[utf8]{inputenc}
% Fix for l3regex.syn error
\usepackage{expl3}
\expandafter\def\csname ver@l3regex.sty\endcsname{}
% Hyperlinks and URLs
\usepackage[hyphens]{url}
\usepackage[colorlinks, breaklinks]{hyperref}
\hypersetup{citecolor=blue, linkcolor=green, urlcolor=magenta, filecolor=cyan, pdfpagemode=UseNone} % Link colors and hide bookmarks
% Citation Style
\usepackage{apacite}
% Math Packages
\usepackage{amsmath, amssymb, amsfonts, amsthm, amsbsy, latexsym}
\usepackage{mathtools} % $mathclap{}$
\DeclareMathOperator*{\plim}{plim} % Custom operator
% Graphics and Figures
\usepackage{graphicx, caption, subcaption} % Include figures \includegraphics[scale=1]{figurename.png}
\usepackage{stackengine, scalerel} % Alternative to \widehat: \stackon
\usepackage{tikz, pgfplots} % TikZ and plotting
\pgfplotsset{compat=newest}
\usetikzlibrary{shapes.misc, positioning, arrows.meta, fillbetween}
% Tables
\usepackage{booktabs, multirow, longtable, array, tabularx, ragged2e} % Enhanced tables
\usepackage{diagbox, makecell} % Diagonal cells and customized table cells
% Lists and Enumerations
\usepackage{enumitem} % Customizable lists with labels: \begin{enumerate}[label=(\roman*)] \begin{enumerate}[label=\fbox{\Roman*}] \begin{enumerate}[label=(\alph*)]
% Text Formatting
\usepackage[dvipsnames]{xcolor} % Custom text colors \textcolor{red}{text} \definecolor{myblue}{RGB}{0, 0, 255}
\usepackage{coloring}
\usepackage{nicefrac} % Inline fractions: \nicefrac{1}{2}
\usepackage{parskip} % Skip lines between paragraphs
\usepackage{indentfirst} % Indent first paragraph after \section{}
\usepackage{upref} % Non-italicized references
\usepackage{cancel} % Strikeout in equations: \cancel{}
% Miscellaneous
\usepackage{verbatim} % Hide blocks: \begin{comment} ... \end{comment}
\usepackage{lineno} % Number lines: \begin{linenumbers} ... \end{linenumbers}
\usepackage[nodayofweek]{datetime} % Short date and time format: \today \currenttime
\usepackage{listings} % Code listings
\usepackage{float} % Control figure placement
\pagenumbering{gobble} % No page numbers
\usepackage[fixed]{fontawesome5}
% I usually forget
	% Write above or below in math mode: $\overset{\mathrm{def}}{=}$ $\underset{\mathrm{def}}{=}$
	% Set column width in tables: \begin{tabular}{|c|p{3cm}|} ... \end{tabular}
	% Scale tables: \scalebox{0.8}{}
	% Multiline expression underneath: $\sum_{\mathclap{\substack{i=0 \\ i \neq 5}}}^{10} a_i$. Use $\displaystyle\sum$ if necessary
	% Avoid all inline math mode: \everymath{\displaystyle}
% Custom commands
\newcolumntype{R}[1]{>{\RaggedRight}p{#1}}
\newcommand{\innerproduct}[2]{\left\langle #1, #2 \right\rangle}
\newcommand{\norm}[1]{\left\lVert #1 \right\rVert}
\newcommand{\normalized}[1]{#1_\circ}
\newcommand{\dg}[1]{\mathrm{dg}\left(#1\right)}
\newcommand{\diag}[1]{\mathrm{diag}\left(#1\right)}
\newcommand{\tr}[1]{\mathrm{tr}\left(#1\right)}
\newcommand{\dimension}[1]{\mathrm{dim}\left(#1\right)}
\newcommand{\col}[1]{\mathrm{col}\left(#1\right)}
\newcommand{\kernel}[1]{\mathrm{ker}\left(#1\right)}
\newcommand{\rk}[1]{\mathrm{rk}\left(#1\right)}
\newcommand{\swp}[1]{\mathrm{SWP}\left(#1\right)}
\newcommand{\etr}[1]{\mathrm{exp}\left(\tr{#1}\right)}
\newcommand{\vecv}[1]{\mathrm{vec}\left(#1\right)}
\newcommand{\vech}[1]{\mathrm{vech}\left(#1\right)}
\newcommand{\red}[1]{\textcolor{red}{#1}}
\newcommand{\green}[1]{\textcolor{green!50!black}{#1}}

\begin{document}

\section*{\faScroll Errata: Bayesian Econometrics - Gary Koop (2003)}

\begin{longtable}{|R{.36\textwidth}|R{.58\textwidth}|} \hline
	\multicolumn{2}{|l|}{Original errata (1\textsuperscript{st} January, 2009): \href{https://epdc.wiley.com/pub/docs/3243ba3d-b457-40a0-ae6b-63d7ae78a1b7}{epdc.wiley.com/pub/docs/3243ba3d-b457-40a0-ae6b-63d7ae78a1b7}} \\
	\multicolumn{2}{|l|}{This errata (\today): \href{https://github.com/zekiakyol/compact-erratas}{github.com/zekiakyol/compact-erratas}} \\ \hline \hline
	\textbf{\faMapMarker* Location} & \textbf{\faMarker Correction} \\ \hline 
	\textbf{Page 1}, Second equation & \red{$p(A,B)=p(B\mid A)p(B)$} should be \green{$p(A,B)=p(B\mid A)p(A)$}. \\ \hline	
	\textbf{Page 5}, 8 lines after Equation (1.8) & The probability  \red{$p(y \mid y^*)$} should be \green{$p(y^*\mid y)$}. \\ \hline	
	\textbf{Page 20}, Equation (2.18) &  \red{$\displaystyle var(h\mid y)=\frac{2\overline{s}^{-2}}{\overline{\nu}}$} should be \green{$\displaystyle var(h\mid y)=\frac{2\overline{s}^{-4}}{\overline{\nu}}$}. \\
	& This error was also made in the MATLAB code used for the empirical example on pages 28--30, so the reported results are very slightly off. \\ \hline	
	\textbf{Pages 23 and 269}, Jeffreys prior & \red{$\displaystyle p(\beta,h)=\frac{1}{h}$} should be  \green{$\displaystyle p(\beta,h)=\frac{1}{\sqrt{h}}$}. \\ \hline	
	\textbf{Page 36}, Equation (3.7) & The \red{$h^{\frac{1}{2}}$} term in the likelihood function should be \green{$h^{\frac{k}{2}}$}. \\ \hline	
	\textbf{Page 37}, Equation (3.19) & \red{$\displaystyle var(h\mid y)=\frac{2\overline{s}^{-2}}{\overline{\nu}}$} should be \green{$\displaystyle var(h\mid y)=\frac{2\overline{s}^{-4}}{\overline{\nu}}$}.  \\ \hline	
	\textbf{Page 38}, Equation (3.24) & \red{$\displaystyle p(\beta,h)\propto \dfrac{1}{h}$} should be \green{$\displaystyle p(\beta,h)\propto h^{\dfrac{k-2}{2}}$}. \\ \hline
	\textbf{Page 42}, Equation (3.37) & \red{$\displaystyle PO_{12}=\frac{(|X_1'X_1|)^{-\frac{1}{2}}(\nu_1s_1^2)^{-\frac{N}{2}}p(M_1)}{(|X_2'X_2|)^{-\frac{1}{2}}(\nu_2s_2^2)^{-\frac{\overline{N}}{2}}p(M_2)}$} should be  \\ 	
	& \green{$\displaystyle PO_{12}=\frac{(|X_1'X_1|)^{-\frac{1}{2}}(\nu_1s_1^2)^{-\frac{N}{2}}p(M_1)}{(|X_2'X_2|)^{-\frac{1}{2}}(\nu_2s_2^2)^{-\frac{N}{2}}p(M_2)}$}. \\ \hline
	\textbf{Page 45}, First line after Equation (3.38) & For notation consistency, \red{$N(0,h^{-1}I_T)$} should be \green{$N(0_T,h^{-1}I_T)$}. \\ \hline	
	\textbf{Page 45}, Equation (3.39) & \red{$S$} should be \green{$T$}. \\ \hline	
	\textbf{Page 46}, Line 1 & The reference to \red{(3.38)} should refer to \green{(3.39)}. \\ \hline	
	\textbf{Page 51}, Table 3.2 & The estimated standard deviations of the prior and posterior $h$ are slightly off. In MATLAB file \texttt{ch3post.m}, line 58, \red{\texttt{hvar=2/v1s12}} should be \green{\texttt{hvar=2/(v1*s12\textasciicircum2)}}. \\ \hline	
	\textbf{Page 60}, 2 lines above Equation (4.3) & The text should refer to \green{``(3.3), (4.1) and (4.2)''}, not \red{``(3.2), (4.1) and (4.2)''}. \\ \hline
	\textbf{Page 66}, Last line & \red{``posterior involves is a mixture''} should be \\
	& \green{``posterior involves a mixture''}. \\ \hline	
	\textbf{Page 66}, Equation (4.14) & \red{$\displaystyle CD=\frac{\widehat{g}_{S_A}-\widehat{g}_{S_C}}{\dfrac{\widehat{\sigma}_A}{\sqrt{S_A}}+\dfrac{\widehat{\sigma}_C}{\sqrt{S_C}}}$} should be \green{$\displaystyle CD=\frac{\widehat{g}_{S_A}-\widehat{g}_{S_C}}{\sqrt{\dfrac{\widehat{\sigma}_A^2}{S_A}+\dfrac{\widehat{\sigma}_C^2}{S_C}}}$}. \\ \hline	
	\textbf{Page 79}, 2 lines above Equation (4.39) & The reference to \red{(4.38)} should refer to \green{(4.37)}. \\ \hline	
	\textbf{Page 84}, Line 4 from top & \red{``weights are then calculated as (4.36)''}  should be \\
	& \green{``weights are then calculated as specified in (4.38).''} \\ \hline	
	\textbf{Page 84}, Line 8 from top &  \red{``as in (4.35), we can $\ldots$''} should  be \green{``as in (4.37), we can $\ldots$''}. \\ \hline	
	\textbf{Page 103}, Line 18 from top & Delete the expectation symbols. The corrected text is \green{``we are setting $g(y)=Skew(y)$ or $Kurt(y)$ and $g(y^\dagger)=Skew(y^\dagger)$ or $Kurt(y^\dagger)$.''} \\ \hline	
	\textbf{Page 103}, Posterior predictive $p$-value discussion & The description is correct as written but may be misleading: the stated procedure calculates the probability of values more extreme than the observed test statistic in one tail only of its posterior distribution. Under the definition used here, the posterior predictive $p$-value lies between $0$ and $0.5$. A more conventional $p$-value can be obtained by doubling it, which should work well when the distribution is approximately symmetric. \\ \hline	
	\textbf{Page 106}, Equation (5.22) & \red{$\displaystyle f(\theta)=\frac{1}{p(2\pi)^{\frac{k}{2}}}|\widehat{\Sigma}|^{-\frac{1}{2}}\exp\left[-\frac{1}{2}(\widehat{\theta}-\theta)'\widehat{\Sigma}^{-1}(\widehat{\theta}-\theta)\right]1(\theta\in\widehat{\Theta})$} should be \\ 
	& \green{$\displaystyle f(\theta)=\frac{1}{(1-p)(2\pi)^{\frac{k}{2}}}|\widehat{\Sigma}|^{-\frac{1}{2}}\exp\left[-\frac{1}{2}(\widehat{\theta}-\theta)'\widehat{\Sigma}^{-1}(\widehat{\theta}-\theta)\right]1(\theta\in\widehat{\Theta})$}. \\ \hline	
	\textbf{Page 107}, Line 3 & Replace with \green{``for the noninformative prior given in (5.4), and$\ldots$''}. \\ \hline	
	\textbf{Page 107}, Paragraph 2, Line 4 from bottom & The reference should be to \green{(5.6)}. \\ \hline	
	\textbf{Page 120}, Equation (6.8) & Insert the missing factor \green{$|\Omega|^{-1/2}$}, and replace \red{$h^{1/2}$} with \green{$h^{k/2}$}. \\ \hline	
	\textbf{Page 129}, Description of Student-$t$ errors & The statement \red{``$\ldots$ linear regression model with independent Student-$t$ errors''} formally holds for fixed $\nu_\lambda$. When $\nu_\lambda$ is treated as an unknown parameter and assigned a hierarchical prior, the resulting distribution is a mixture of Student-$t$ distributions. \\ \hline	
	\textbf{Page 153}, Second equation from bottom &  \red{$\displaystyle \overline{\beta}=\overline{V} \left(\underline{V}_{\beta}^{-1}\underline{\beta}+h\sum_{i=1}^{N}\widetilde{X}_i'[y_i-\alpha_i\iota_T]\right)$} should be \\
	& \green{$\displaystyle \overline{\beta}=\overline{V}_{\beta}\left(\underline{V}_{\beta}^{-1}\underline{\beta}+h\sum_{i=1}^{N}\widetilde{X}_i'[y_i-\alpha_i\iota_T]\right)$}. \\ \hline	
	\textbf{Page 160}, Equation (7.36) & Replace with \green{$\displaystyle \widehat{p}(\theta_2^*\mid y,\theta_1^*)=\frac{1}{S^*}\sum_{s^*=1}^{S^*}p\!\left(\theta_2^*\mid y,z^{(s^*)},\theta_1^*\right)$}. \\ \hline
	\textbf{Page 183}, Third line from bottom & Use \green{$E(\alpha_t\mid\alpha_{t-1})=\alpha_{t-1}$}. \\ \hline	
	\textbf{Page 184}, Line 5 & \red{$u_t$} should be \green{$\varepsilon_t$}. \\ \hline	
	\textbf{Page 190}, 15 lines from bottom & The reference to \red{(8.10)} should be \green{(8.19)}. \\ \hline	
	\textbf{Page 198}, Equation (8.42) & Replace with \green{$\displaystyle K_t=(T_tP_tZ_t'+J_tG_t')D_t^{-1}$}. \\ \hline	
	\textbf{Page 199}, Equation (8.44) & Replace with \green{$\displaystyle P_{t+1}=T_tP_t(T_t-K_tZ_t)'+J_t(J_t-K_tG_t)'$}. \\ \hline	
	\textbf{Page 199}, Equation (8.47) & Replace with \green{$\displaystyle V_t=F_t\left(G_t'D_t^{-1}Z_t+[J_t-K_tG_t]'U_t[T_t-K_tZ_t]\right)$}. \\ \hline	
	\textbf{Page 202}, Empirical example & A small error in the MATLAB code used for this empirical example makes \red{Table 8.1 and Figure 8.3} incorrect; see the computer-code errata below. \\ \hline	
	\textbf{Page 242}, Equation for $p(\eta\mid y)$ at bottom & The equation is incorrect. Use the general form given in \green{(8.21) on page 189}. \\ \hline	
	\textbf{Page 268}, 9 lines from bottom & \red{$(y_1,\ldots,y_T)'$} should be \green{$(y_1,\ldots,y_N)'$}. \\ \hline	
	\textbf{Page 271}, Below Equation (11.9) & The degrees of freedom should be \green{$N-1$}, not \red{$N$}. \\ \hline	
	\textbf{Page 271}, Equation (11.12) & The \red{$\iota_T$'s} should be \green{$\iota_N$'s}. \\ \hline	
	\textbf{Page 271}, Between Equations (11.11) and (11.12) & Replace the formula for $P_{X_r}$ with
	\green{$\displaystyle P_{X_r}=M_1-X_r(X_r'X_r)^{-1}X_r'$}, where \green{$\displaystyle M_1=I_N-\frac{\iota_N\iota_N'}{N}$}. \\ \hline	
	\textbf{Page 290}, Information matrix & The information matrix is based on the \green{log likelihood}, not the likelihood itself. Two equations on this page should therefore contain log likelihoods instead of likelihoods. \\ \hline	
	\textbf{Page 320}, Theorem B.2 & Replace the theorem with: \green{Let $X$ and $Y$ be two random variables, $g(\cdot)$ and $h(\cdot)$ be two functions, and $a$ and $b$ be constants. Then $E[ag(X)+bh(Y)]=aE[g(X)]+bE[h(Y)]$, and $var[ag(X)+bh(Y)]=a^2 var[g(X)]+b^2 var[h(Y)]$ if $X$ and $Y$ are independent.} \\ \hline	
	\textbf{Page 322}, Definitions B.13 and B.14 & Definitions B.13 and B.14 repeat the definitions on page 321. \\ \hline	
	\textbf{Page 327}, Multivariate Normal p.d.f. & \red{$2\pi^{k/2}$} should be \green{$(2\pi)^{k/2}$}. \\ \hline	
	\textbf{Page 329}, Wishart distribution & The Wishart distribution is only defined for \green{$\nu\geq N$}. \\ \hline	
	\textbf{Page 339}, Hobert and Casella reference & The reference is in volume \green{91}, not \red{96}. \\ \hline	
	\textbf{Page 344}, Zellner and Tobias reference & Replace with \green{Zellner, A. and Tobias, J. (2001). ``Further Results on Bayesian Method of Moments Analysis of the Multiple Regression Model,'' \textit{International Economic Review}, 42, 121--139.} \\ \hline  \hline
	\multicolumn{2}{|l|}{\textbf{Errors in the Computer Code on the Book's Website}} \\ \hline	
	\textbf{\texttt{chapter8c.m}} & The $\lambda$'s are not passed correctly into the \texttt{djs.m} algorithm. This causes the empirical results on page 202 to be based on incorrect code, so those results are incorrect. \\ \hline	
	\textbf{\texttt{chapter7d.m}}, Lines 67 and 148 & In the Chib marginal-likelihood calculation, the log determinants of the prior and posterior covariance matrices are multiplied by \red{\texttt{kx}}; this multiplier should \green{not be present}. Koop notes that the fix had been made, but as of January 2009 the website version had not been updated. \\ \hline	
	\textbf{\texttt{chapter7c.m}} & A reader reports an error in the drawing of the slope coefficients. \\ \hline	
	\textbf{\texttt{chapter7d.m}}, \texttt{sigterm1} & There is a minor error in the step defining \texttt{sigterm1} in two places. See \green{equation (6.54)} for the correct formulation. \\ \hline	
\end{longtable}

\end{document}